\section{Discussion}

We presented \olmix, addressing two key challenges in data mixing for language models: configuring effective mixing methods and efficiently updating mixtures as domain sets evolve. We conduct a large empirical study of design choices and introduce mixture reuse. 

\textbf{Limitations.} First, our work focuses exclusively on the offline mixing schema. While this schema is widely adopted, other approaches, such as those that use a dynamic mix to explore the mixture weight space~\citep{xie2023doremioptimizingdatamixtures, fan2024dogedomainreweightinggeneralization}, may benefit from different design choices. Second, our theoretical analysis of mixture reuse assumes log-linear regression models. The analysis should extend naturally to other parametric models like AutoScale~\citep{kang2025autoscalescaleawaredatamixing} and BiMix~\citep{ge2025bimixbivariatedatamixing}, but it is less clear how to extend it to non-parametric models that yield non-convex and non-differentiable mixing objectives.
Third, \partialmixreuse requires determining the subset of unaffected domains to recompute. Some cases are quite intuitive; for instance, recomputing DCLM:software\_development when Stack-Edu code data is added. However, we do not provide a general automated approach for determining what to recompute, limiting \partialmixreuse's applicability without domain expertise. 

\textbf{Future work.} First, we aim to extend mixture reuse and our design choice study to online mixing methods that adjust mixtures during training. This is in contrast with this work, which focuses on the LM development that occurs \textit{before} the final model starts training. Second, we envision co-design of data mixing with other LM development workflows, such as using mixture performance as feedback for data quality filtering or domain discovery. Third, validating our findings at larger model scales would improve confidence in the reliability of our design choices and reuse mechanisms.


\section*{Author Contributions}

Mayee F. Chen led the project during an internship at Ai2: she developed the theoretical framework, implemented all methods, designed and ran all experiments, validated the approach at scale for OLMo 3, and wrote the paper. Kyle Lo and Luca Soldaini were the primary mentors and worked closely with Mayee on developing the methodology, designing experiments, integrating into Olmo 3, and writing the paper. Tyler Murray built and maintained the training infrastructure that enabled large-scale experimentation. David Heineman designed small-scale evaluation metrics for data mixing. Matt Jordan contributed to mixing methodology. Hannaneh Hajishirzi and Christopher R\'{e} provided research guidance and helped frame the work.

\section*{Acknowledgments}

We thank Neel Guha, Junmiao Hu, Ronny Junkins, Pang Wei Koh, Jerry Liu, Roberto Garcia Torres, Alex Wettig, Steven Euijong Whang, and Michael Zhang for helpful feedback and discussions. 
We thank the Olmo 3 team, particularly Allyson Ettinger, for iterative feedback and experimental validation of Olmix.
MC was supported in part by the Office of Naval Research (ONR) under No. N000142312633 (Deep Signal Processing). Any opinions, findings, and conclusions or recommendations expressed in this material are those of the authors and do not necessarily reflect the views, policies, or endorsements, either expressed or implied, of ONR or the U.S. Government.